\section{The Evaluation Framework}
\label{sec:framework}

\subsection{Measurement formalization}

Let $M$ be a set of models, $T$ a set of tasks, and $C$ a space of
evaluation configurations, where a configuration $c \in C$ fixes the prompt
variant, few-shot count, chat template, inference backend, quantization,
decoding strategy, seed, and hardware/software environment. For a benchmark
slice of $n$ items with gold answers, the measured score of model $m$ on
task $t$ under configuration $c$ is
\begin{equation}
S(m, t, c) \;=\; \frac{1}{n_{t}} \sum_{i \in D_t} \mathbb{1}\!\left[
\hat{y}_i(m, c) = y_i \right],
\end{equation}
where $D_t$ is the frozen item set for task $t$ ($|D_t| = 200$),
$\hat{y}_i(m,c)$ is the extracted answer from the generated text, and $y_i$
is the gold. The overall score pools the four tasks ($n = 800$). Conventional
reporting emphasizes a single value $S(m, t, c_0)$ under one fixed
configuration $c_0$; this work measures the set
$\{S(m,t,c) : c \in C_{\text{measured}}\}$ over a controlled neighborhood of
$c_0$ and asks what it implies for scores and rankings.

Define the pooled score $\bar S(m,c)=|T|^{-1}\sum_{t\in T}S(m,t,c)$;
this equals item-pooled accuracy because the task slices have equal size.
A model \emph{ranking} under configuration $c$ is $R_c(m)$, the competition
rank (1 = best) by $\bar S(m,c)$. Equal scores receive equal ranks; subsequent
ranks skip the tied positions. Ranking stability across configurations is
summarized by Kendall's $\taub$ between the control ranking $R_{c_0}$ and
each variant ranking $R_c$ \citep{kendall1938new}, together with the number
of pairwise reversals. We additionally report the pairwise win rate
\begin{equation}
P(A > B) \;=\; \frac{\left|\{c \in C_{\text{shared}} : \bar S(A,c) > \bar S(B,c)\}\right|}{|C_{\text{shared}}|},
\end{equation}
over configurations where both models were measured.

\subsection{Interval-aware comparison}

Every accuracy is reported with a Wilson score interval
\citep{wilson1927probable}, which remains well behaved at the small $n$ and
extreme proportions that occur in our cells. Marginal interval overlap is
reported descriptively, not as a test of equal model performance. Overlap
neither establishes equivalence nor rules out a paired difference. Wilson
intervals assume binomial sampling; the fixed heterogeneous task mixture and
possible dependence between items limit this interpretation. Paired
item-level comparisons use McNemar's test \citep{mcnemar1947note} with
continuity correction on the discordant pairs, which tests the paired
hypothesis that the two configurations produce the same per-item correctness
probability; because it uses pairing, it can detect changes even when
unpaired intervals overlap. Effect sizes are reported as Cohen's $h$
\citep{cohen1988statistical}. Because we evaluate a family of 27 OFAT
contrasts against controls, we correct for multiple testing with
Holm--Bonferroni \citep{holm1979simple} and report Benjamini--Hochberg
false-discovery-rate-adjusted values \citep{benjamini1995controlling} in the
appendix. Ranking-level uncertainty is quantified with a seeded bootstrap
that resamples configurations (not items) and recomputes rankings.

\subsection{What the framework records}

Every run writes a JSON artifact containing: the full settings block (model
ID, tokenizer ID, \texttt{revision}, backend, chat template, prompt variant,
system prompt, \texttt{max\_new\_tokens}, seed, dtype, quantization,
temperature, top-$p$, batch size, data path, and experiment ID); package
versions (PyTorch, Transformers, vLLM, harness); the hardware manifest
(platform, Python version, GPU model, CUDA availability); per-item records
(item ID, task, gold, predicted, correct, raw generation, and per-item
timing); per-task and overall accuracies with Wilson intervals; and a
latency block. A 16-hex configuration hash is derived from the settings
fields that affect the measurement; the registry stores one row per
completed configuration, deduplicated by hash, and each row points at its
artifact. Every artifact also embeds an \emph{incomparability} statement
listing the dimensions that would make two scores non-comparable. These
records are what allow every number in this paper to be traced to its
configuration (\cref{sec:repro}).

\subsection{Planned nested factorial experiment}
\label{sec:factorial-plan}

OFAT estimates simple effects at the control settings, not marginal effects
averaged over other factors. The existing follow-up YAML specifies three
prompts (default, concise, five-shot), three runtime bundles (HF
unquantized, HF INT8, vLLM unquantized), and three greedy seeds (0, 1, 2):
$3\times3\times3=27$ cells per model, or 81 for all three models.
The runtime factor is bundled because vLLM quantization is unsupported in
this harness. Prompt $\times$ backend is estimable at no quantization;
Prompt $\times$ quantization is estimable within HF. Backend $\times$
quantization cannot be identified from this support pattern.

The plan in \texttt{configs/experiments/t4\_factorial.yaml} uses a separate
200-item slice (100 ARC-Easy and 100 GSM8K), not the historical 800-item
suite. Greedy seeds are determinism checks, not independent stochastic
replicates. Sampling would require a separately declared decoding arm.
Before execution, model revisions, software, hardware and scoring must be
frozen. Interaction analysis should compare paired differences-in-differences
on identical item IDs, reporting uncertainty with item-level pairing and
accounting for repeated responses and related items; seed duplicates must
not inflate the effective sample size. No factorial inference was run in
this revision: the available machine has an Apple GPU, not the required
CUDA/T4 environment. These remain prospective analyses, not results.


\section{Related Work}
\label{sec:related}

\paragraph{Evaluation harnesses and holistic benchmarks.}
The lm-evaluation-harness popularized configurable, task-level LLM
evaluation \citep{gao2021framework}, and its maintainers document evaluation
pitfalls including prompt-format drift, backend differences, and missing
error bars \citep{biderman2024lessons}. HELM frames evaluation as a matrix
of scenarios and metrics \citep{liang2022holistic}, BIG-bench scales task
diversity \citep{srivastava2022beyond}, and Dynabench argues for dynamic,
documented benchmark iteration \citep{kiela2021dynabench}. These systems provide configurable evaluation protocols. Our contribution
is a bounded joint audit of operator choices and stored item outcomes,
not a claim that existing harnesses cannot configure these dimensions.

\paragraph{Prompt sensitivity.}
\citet{sclar2024quantifying} measure performance spread under
meaning-preserving prompt-format changes and find differences of up to tens
of accuracy points, recommending ranges rather than single numbers.
\citet{mizrahi2024state} evaluate 20 LLMs across 39 tasks under diverse
prompts and show that single-prompt rankings can be brittle.
\citet{zhu2024promptrobust} show robustness failures under adversarial
prompt perturbations. Our study differs in scope---we use three deliberately
distinct instruction configurations rather than large format
distributions---but shares the premise that the prompt is part of the
measurement, and extends the analysis to paired per-item artifacts,
corrected paired tests, and explicit ranking-stability statistics on frozen
slices.

\paragraph{Statistical treatment of evaluations.}
\citet{miller2024adding} formalizes evaluations as experiments and
recommends super-population confidence intervals for reported scores;
tinyBenchmarks \citep{polo2024tinybenchmarks} studies how few items suffice
to estimate benchmark performance. We adopt interval reporting and add
paired item-level tests (McNemar) with family-wise correction, effect
sizes, and a config-level bootstrap for ranking stability.

\paragraph{Backends and quantization.}
vLLM's PagedAttention serving stack \citep{kwon2023efficient} and
bitsandbytes INT8/INT4 quantization \citep{dettmers2022llm,dettmers2023case}
are widely used in open-weight evaluation; k-bit scaling laws quantify
aggregate zero-shot accuracy trade-offs across many models. Our study
measures these dimensions as factors on fixed slices with paired statistics,
including the observation that our harness deliberately calls vLLM in
completion mode so that both backends consume byte-identical rendered
prompts.

\paragraph{Sampling temperature.}
\citet{renze2024temperature} find no statistically significant effect of
temperature on MCQA problem solving between 0.0 and 1.0 at their scale. Our
sampling arm complements this with per-item churn measurements: net drift
can be small while 9.9--24.1\% of items change correctness state.


\paragraph{Reproducibility and benchmark reliability.}
\citet{biderman2024lessons} document sensitivity to evaluation setup and
recommend transparent reporting of evaluation conventions. Our paired
artifact audit operationalizes this recommendation, while distinguishing
stored-output reproducibility from historical inference replication.
LiveBench \citep{white2024livebench}, first released in 2024 and revised in
2025, uses frequently updated questions and objective ground-truth scoring
to limit contamination and judging bias. This addresses a different threat
from configuration sensitivity: our frozen public slices support paired
comparisons but do not establish freedom from contamination.

\paragraph{Quantization beyond the measured implementation.}
AWQ \citep{lin2023awq} uses activation information to protect salient weight
channels in low-bit weight-only quantization and evaluates language-model
and domain-specific benchmarks. Alongside k-bit scaling laws
\citep{dettmers2023case}, this shows why precision results must be tied to a
specific method and workload. We test bitsandbytes, not AWQ, and make no
claim that the small corrected effects here generalize to other quantizers.
